% chapters/04_methodik.tex
%
% Methodikkapitel. Angaben zu Implementierung, Konfiguration und
% Laufumgebung sind gegen Quellcode und gespeicherte Laufmetadaten
% abgeglichen; die Zuordnung von Läufen zu Commits, Revisionen und
% Artefakten steht in anhang.tex (Label app:reproduzierbarkeit).
% Benötigte Pakete in der Präambel: xcolor, tabularx, listings, booktabs.

\chapter{Methodik}
\label{chap:methodik}

Anhang~\ref{app:reproduzierbarkeit} ordnet jedem hier beschriebenen Lauf
Konfiguration, Codestand,
Revisionen und Ausgabeverzeichnis zu;
Anhang~\ref{app:durchfuehrung} ergänzt Durchführungsdetails.

\section{Methodischer Ansatz}
\label{sec:methodischer-ansatz}

Die Arbeit verfolgt einen empirisch-vergleichenden Ansatz: Statt neue
Architekturen oder Trainingsverfahren zu entwickeln, stellt sie
etablierte NER-Verfahren unter möglichst einheitlichen Bedingungen
gegenüber.
Tragfähige Vergleiche setzen voraus, dass Datensatz, Datenaufteilung und
Evaluationsprotokoll über alle Konfigurationen konstant bleiben.

\paragraph{Drei Vergleichsregime.}
\begin{enumerate}
    \item \textbf{Encoder-Fine-Tuning:} Vortrainierte Transformer-Encoder
    werden zusammen mit einem linearen Klassifikationskopf vollständig
    feinjustiert und liefern direkt IOB2-Tag-Sequenzen
    \cite{devlin2019bert,keraghel2024recent}.

    \item \textbf{LLM-Zero-Shot:} Veröffentlichte, bereits vor- und
    nachtrainierte Decoder-LLMs werden ohne aufgabenspezifisches Training
    und ohne Demonstrationsbeispiele im Prompt eingesetzt und geben die
    Entitäten als JSON aus.
    Das Regime beschreibt die Leistung dieser Checkpoints unter den
    gewählten Prompt- und Decoding-Bedingungen
    \cite{brown2020language,zhao2023survey}.

    \item \textbf{LLM-Fine-Tuning mit LoRA (QLoRA-Regime):} Dieselben LLMs werden über
    trainierbare Low-Rank-Adapter angepasst, während die ursprünglichen
    Gewichte eingefroren bleiben \cite{hu2021lora}.
    Wegen der Modellgröße wird das Regime als QLoRA mit
    4-Bit-quantisiertem Basismodell umgesetzt \cite{dettmers2023qlora}
    (Abschnitt~\ref{sec:lora-konfiguration}) und im Folgenden
    QLoRA-Regime genannt.
\end{enumerate}

Encoder werden ausschließlich feinjustiert, da Token-Klassifikatoren mit
neu initialisiertem Kopf ohne Training keine NER-Vorhersagen erzeugen und
somit kein Zero-Shot-Pendant besitzen; Generalisten wie GLiNER
(Abschnitt~\ref{sec:encoder_ner_related}) bleiben als eigene Modellklasse außerhalb des
Vergleichs.

\paragraph{Experimentmatrix.}
Aus den Regimen ergeben sich acht Konfigurationen
(Tabelle~\ref{tab:experiment-matrix}).
Die fünf trainierten Konfigurationen werden mit den Seeds 42, 123 und 456
jeweils unabhängig trainiert und evaluiert (15 Trainingsläufe); die drei
Zero-Shot-Konfigurationen wurden je einmal mit Seed~42 ausgewertet.
Die Hauptauswertung umfasst damit 18 Testevaluationen.

Der Hauptstudie gingen explorative Vorläufe der fünf trainierten
Konfigurationen mit Seed~42 voraus, die der Stabilisierung der Pipeline
dienten und nicht in die Auswertung eingehen
(Tabelle~\ref{tab:datennutzung}, Tabelle~\ref{tab:app-vorlaeufe}).
Die Zero-Shot-Läufe stammen aus einem eindeutig identifizierbaren Commit
ohne lokale Änderungen (Anhang~\ref{app:reproduzierbarkeit}) und sind Teil
der Hauptauswertung.

\begin{table}[ht]
    \centering
    \small
    \caption[Untersuchte Konfigurationen]{Untersuchte Konfigurationen.
    Größe: relative Größenklasse innerhalb der Modellfamilie
    (Abschnitt~\ref{sec:groessenklassen}).
    Parameter: Textmodell ohne Adapter, bei Qwen3.5 ohne den nicht
    geladenen Vision-Encoder (Abschnitt~\ref{sec:llm-auswahl}).
    Trainierbar: bei LoRA die Adapterparameter.
    NF4: 4-Bit-quantisiertes Basismodell; die Adapter werden mit
    BF16-Mixed-Precision trainiert (Abschnitt~\ref{sec:lora-konfiguration}).}
    \label{tab:experiment-matrix}
    \setlength{\tabcolsep}{4pt}
    \begin{tabular}{clllrrr}
        \toprule
        \textbf{Größe} & \textbf{Modell} & \textbf{Regime} &
        \textbf{Präzision} & \textbf{Parameter} & \textbf{Trainierbar} &
        \textbf{Läufe} \\
        \midrule
        S & DeBERTa-v3-base  & Fine-Tuning & FP32 & 184\,M   & 184\,M   & 3 \\
        M & DeBERTa-v3-large & Fine-Tuning & FP32 & 434\,M   & 434\,M   & 3 \\
        S & Qwen3.5-0.8B     & Zero-Shot   & BF16 & 0,75\,B  & --       & 1 \\
        M & Qwen3.5-4B       & Zero-Shot   & BF16 & 4,21\,B  & --       & 1 \\
        L & Qwen3.5-27B      & Zero-Shot   & BF16 & 26,90\,B & --       & 1 \\
        S & Qwen3.5-0.8B     & QLoRA       & NF4  & 0,75\,B  & 6,4\,M   & 3 \\
        M & Qwen3.5-4B       & QLoRA       & NF4  & 4,21\,B  & 21,2\,M  & 3 \\
        L & Qwen3.5-27B      & QLoRA       & NF4  & 26,90\,B & 159,4\,M & 3 \\
        \bottomrule
    \end{tabular}
\end{table}

\paragraph{Variablen, Konstanten und Störgrößen.}
Die Konfigurationen unterscheiden sich in \emph{Modellfamilie und
Architektur} (Encoder-only gegenüber Decoder-only), \emph{Parameteranzahl}
und \emph{Anpassungsregime}; diese Merkmale werden nicht unabhängig
voneinander variiert. Der Seed dient bei den trainierten
Konfigurationen als Replikationsfaktor.
Abhängige Variablen sind
(1)~strikte Precision, Recall und F$_1$ auf Entitätsebene,
(2)~Trainingsdauer, Inferenzlatenz und GPU-Spitzenspeicher der Inferenz,
(3)~die Zahl gesamter und trainierbarer Parameter sowie
(4)~Kennzahlen zur Formtreue der LLM-Ausgaben.

Konstant sind die Datenbasis (der englische Teil von MultiNERD mit
seinen offiziellen Splits \cite{tedeschi2022multinerd}), das
Evaluationsprotokoll mit derselben Metrikimplementierung für beide
Pipelines (Abschnitt~\ref{sec:performanzmetriken}) und die Versionen der
zentralen Bibliotheken (Abschnitt~\ref{sec:rechenumgebung}).
Alle Konfigurationen werten den vollständigen Testsplit aus.
Nicht konstant sind dagegen der GPU-Typ, die je Modellgröße gewählten
Hyperparameter sowie Prompt, Decoding und numerische Präzision der
Zero-Shot- gegenüber den QLoRA-Läufen.
Abschnitt~\ref{sec:validitaet} fasst diese Abweichungen zusammen.

\paragraph{Aufbau der Pipelines.}
Die Encoder-Pipeline behandelt NER als Token-Klassifikation: Ein
DeBERTa-v3-Encoder mit linearem Klassifikationskopf sagt für jedes Wort
ein IOB2-Tag vorher
\cite{he2021debertav3,wolf2020transformers,devlin2019bert}.
Die LLM-Pipeline behandelt NER als strukturierte Generierung: Das Modell
erhält einen Prompt mit nummerierter Tokenliste und erzeugt ein
JSON-Array von Entitäten, das eine regelbasierte, deterministische
Nachverarbeitung in eine IOB2-Sequenz überführt
\cite{dagdelen2024structuredIE,semin2026spanLabeling}.
Erst diese Rückkonvertierung erlaubt die Evaluation beider Pipelines auf
derselben Repräsentation.
Der Gesamt-F$_1$ bewertet damit die vollständige Vorhersagepipeline
einschließlich Ausgabeformat und Validierung.

\section{Auswahl der Modelle}
\label{sec:modellauswahl}

\subsection{Encoder-Modelle (DeBERTa-v3)}
\label{sec:encoder-auswahl}

Auf der Encoder-Seite werden DeBERTa-v3-base und DeBERTa-v3-large
eingesetzt.
DeBERTa-Varianten zählen auf mehreren NER-Benchmarks zu den
leistungsfähigsten Encodern und dienen häufig als starke Baselines
\cite{keraghel2024recent}.
Base- und Large-Variante erlauben einen Größenvergleich innerhalb
derselben Architektur; einschließlich des Klassifikationskopfes umfassen
sie 183\,855\,391 bzw. 434\,043\,935 Parameter.
Beide Modelle stehen laut ihren Modellkarten unter der MIT-Lizenz
\cite{he2021debertav3card,he2021debertav3cardlarge} und sind über die
\emph{Transformers}-Bibliothek verfügbar \cite{wolf2020transformers}.

\subsection{LLM-Modelle (Qwen3.5)}
\label{sec:llm-auswahl}

Auf der LLM-Seite werden drei Größen der Qwen3.5-Familie verwendet:
Qwen3.5-0.8B, Qwen3.5-4B und Qwen3.5-27B.
Qwen3.5 baut auf der Architektur- und Trainingstradition von Qwen2.5 und
Qwen3 auf \cite{qwen2024qwen25,qwen2025qwen3,qwen2026qwen35blog}; die
Modelle sind laut ihren Modellkarten als offene Gewichte unter der
Apache-2.0-Lizenz verfügbar
\cite{qwen2026card08b,qwen2026card4b,qwen2026card27b}.
Für die Familie sprachen neben der Lizenz die fein abgestufte Größenleiter mit
gemeinsamem Tokenizer und Chat-Template, sodass Größeneffekte nicht mit
Architekturwechseln vermischt werden.
Die drei Stufen wurden so gewählt, dass sie den in der Literatur
untersuchten Bereich von 3 bis 9~Mrd. Parametern
(Abschnitt~\ref{sec:research_gap}) nach unten und oben verlassen und das
Training mit einer einzelnen GPU je Lauf möglich bleibt.
Verwendet werden die veröffentlichten, laut Modellkarten vor- und
nachtrainierten Checkpoints
\cite{qwen2026card08b,qwen2026card4b,qwen2026card27b}; der Begriff
\emph{Basismodell} bezeichnet im Folgenden lediglich die eingefrorenen
Ausgangsgewichte, auf die beim LoRA-Training Adapter aufgesetzt werden.

Laut den Konfigurationsdateien der verwendeten Revisionen sind alle drei
Größen dichte Decoder mit hybrider Attention
(Tabelle~\ref{tab:app-qwen-architektur}): In jedem Block aus vier Schichten
nutzen drei Schichten lineare Attention in Form von Gated DeltaNet, die
vierte vollständige Softmax-Attention mit Grouped-Query-Attention und
Ausgabe-Gating \cite{qwen2026card08b,qwen2026card4b,qwen2026card27b}.

Von den multimodal angelegten Checkpoints laden beide LLM-Regime (Zero-Shot
und QLoRA) mit der Klasse \texttt{AutoModelForCausalLM} nur das Textmodell; Vision-Encoder
und Module für Multi-Token-Prediction werden verworfen
(Tabelle~\ref{tab:app-qwen-architektur} in Anhang~\ref{app:durchfuehrung}).
Alle Parameterangaben beziehen sich daher auf das Textmodell, dessen Größe
von den nominellen Bezeichnungen abweicht
(Tabelle~\ref{tab:experiment-matrix}); die Stufen unterscheiden sich um
den Faktor 5,6 bzw. 6,4.

\subsection{Größenklassen}
\label{sec:groessenklassen}

Zur Strukturierung der Auswertung werden die Modelle in die Klassen
\emph{Small} (S), \emph{Medium} (M) und \emph{Large} (L) eingeteilt
(Tabelle~\ref{tab:experiment-matrix}).
Die Klassen bezeichnen die relative Position innerhalb der jeweiligen
Familie, nicht absolute Parametergrenzen; L ist nur auf der LLM-Seite
besetzt, und die Einteilung ist eine Festlegung dieser Arbeit für die
Analyse in Abschnitt~\ref{sec:einfluss-parameteranzahl}.

\section{Auswahl und Aufbereitung der Datensätze}
\label{sec:datensaetze}

\subsection{Auswahlkriterien}
\label{sec:datenauswahl}

Für MultiNERD sprechen vier Kriterien.
Erstens umfasst das Inventar 15 Entitätstypen
(Tabelle~\ref{tab:datensplits}) und geht damit über Person, Ort und
Organisation hinaus, sodass Unterschiede zwischen den Modellklassen auch
bei seltenen oder semantisch komplexen Typen sichtbar werden können.
Zweitens kombiniert der Datensatz mit Wikipedia- und Wikinews-Artikeln
zwei Textgenres \cite{tedeschi2022multinerd}.
Drittens stehen offizielle \texttt{train}-, \texttt{validation}- und
\texttt{test}-Splits zur Verfügung \cite{tedeschi2022multinerd}, die
künftigen Arbeiten einen festen Bezugspunkt bieten.
Veröffentlichte MultiNERD-Werte beruhen dagegen auf eigener Aufbereitung
oder abweichendem Testumfang; Werte auf den offiziellen Splits sind
nicht auffindbar (Abschnitte~\ref{sec:k3-multinerd}
und~\ref{sec:research_gap}).
Viertens ist der englische Teil mit rund 164\,000 Sätzen groß genug für
vollständiges Fine-Tuning der Encoder ebenso wie für LoRA-Training der
LLMs.
Der Datensatz ist laut Datensatzkarte unter CC~BY-NC-SA~4.0
veröffentlicht und auf nichtkommerzielle Forschung beschränkt
\cite{babelscape2023multinerdcard}.

MultiNERD ist automatisch annotiert \cite{tedeschi2022multinerd}; auch
die Referenz des Testsplits ist somit eine Silberreferenz.
Die von den Autoren manuell geprüfte Stichprobe von 1\,000 zufällig
gezogenen englischen Sätzen wird nicht verwendet.
Sie ist in den veröffentlichten Daten nicht enthalten, und aus welchem
Split ihre Sätze stammen, ist nicht angegeben, sodass Überschneidungen
mit den Trainingsdaten nicht auszuschließen sind.
Typbezogene Aussagen wären auf ihr zudem sehr unsicher
\cite[Tab.~5]{tedeschi2022multinerd}.
Nach dieser Stichprobe unterscheidet sich die Annotationsqualität stark
zwischen den Typen (Abschnitt~\ref{sec:datasets_corpora}).
Die trainierten Konfigurationen nutzen dieselben annotierten
Trainingssätze, alle Systeme werden gegen dieselbe Referenz bewertet;
die Zero-Shot-Konfigurationen erhalten kein Training.
Verlustfunktion, Anpassungsumfang (vollständiges Fine-Tuning gegenüber
Adaptern) und Checkpoint-Auswahl unterscheiden sich jedoch, und die
gemeinsame Silberreferenz garantiert nicht, dass ihre Fehler für beide
Modellfamilien neutral wirken; absolute Werte messen zudem die
Übereinstimmung mit einer automatisch erzeugten Referenz.

\subsection{Datenvorverarbeitung}
\label{sec:datenvorverarbeitung}

Vorverarbeitet wird nur, was die technische Integration erfordert; eine
Textbereinigung oder Normalisierung auf Datensatzebene findet nicht
statt.
Der Datensatz wird über die \emph{Datasets}-Bibliothek in einer festen
Revision (Tabelle~\ref{tab:app-revisionen}) aus dem
Hugging-Face-Repository \texttt{Babelscape/\allowbreak multinerd} geladen,
über das Sprachfeld auf \texttt{en} gefiltert und auf vorhandene Splits
und Spalten sowie exakte Splitgrößen geprüft.
Eine zeilenweise Prüfung beim Laden unterbleibt; die nachträgliche
Prüfung der Rohdaten ergab weder leere Sätze noch Längeninkonsistenzen
oder Label-IDs außerhalb des Inventars.
Die Label-IDs werden über eine feste Liste (\texttt{O}, \texttt{B-PER},
\texttt{I-PER}, \dots, \texttt{I-VEHI}) in der Typreihenfolge von
Tabelle~\ref{tab:datensplits} auf Labelnamen abgebildet.

\paragraph{Encoder-Pipeline.}
Die wortweise vorliegenden Sätze werden mit dem Tokenizer von DeBERTa-v3
in Subwörter zerlegt.
Sequenzen werden auf 256 Subwort-Tokens einschließlich Sondertokens
gekürzt.
Beim Label-Alignment erhält das erste Subwort jedes Wortes dessen Label;
Folge-Subwörter und die Sondertokens \texttt{[CLS]} und \texttt{[SEP]}
erhalten den Wert $-100$ und werden in Verlust und Evaluation ignoriert
\cite{devlin2019bert,keraghel2024recent}.
Zur Evaluation werden die Vorhersagen der ersten Subwörter auf die
Wortebene übertragen, sodass Vorhersage und Referenz genau ein Label
pro Wort umfassen.
Wörter jenseits dieser Längengrenze entfallen; im Testsplit betrifft
dies keinen Satz, sodass alle 347\,025 Testwörter evaluiert werden.
Für Trainings- und Validierungssplit wurde die Zahl der gekürzten Sätze
nicht protokolliert.

\paragraph{LLM-Pipeline.}
Jedes Beispiel wird in eine Chat-Struktur aus Systemnachricht (Aufgabe
und Typinventar) und Nutzernachricht (nullbasiert nummerierte Tokenliste)
überführt, im Training ergänzt um eine Assistentennachricht mit der
Zielantwort (Abschnitt~\ref{sec:prompt-design}); die Nummerierung der
Eingabetokens nach \textcite[Tab.~9]{ramponi2025fallacy} erlaubt die
Rückabbildung generierter Spans auf die Tokenfolge
\cite{semin2026spanLabeling}.
Die Zielantwort wird deterministisch aus der IOB2-Sequenz abgeleitet: Jede
Entität wird in Satzreihenfolge mit den vier Feldern aus
Abschnitt~\ref{sec:prompt-design} in ein JSON-Array geschrieben; Sätze
ohne Entitäten erhalten das leere Array.
Für die Seed-42-Läufe aller trainierten Konfigurationen wurde bestätigt,
dass die verwendeten Referenzsequenzen für sämtliche Testsätze exakt der
Originalannotation entsprechen, in der LLM-Pipeline einschließlich der
Hin- und Rückkonvertierung.

\subsection{Datenaufteilung}
\label{sec:datenaufteilung}

Die offiziellen Splits werden unverändert und ohne Teilstichproben
übernommen; die Datenaufteilung ist damit für alle Konfigurationen
identisch und mit Arbeiten vergleichbar, die dieselben Splits verwenden
\cite{tedeschi2022multinerd}.
Die Typverteilung ist stark unausgewogen (Tabelle~\ref{tab:datensplits}):
PER, ORG und LOC stellen 74\,\% der Trainings- und 80\,\% der
Testentitäten, BIO, INST, MYTH und VEHI jeweils weniger als 40
Testentitäten.
Zudem verschiebt sich die Verteilung: Der Anteil von LOC steigt vom
Trainings- zum Testsplit von 29,1\,\% auf 46,9\,\%, der von PER sinkt
von 31,2\,\% auf 20,5\,\%.

\begin{table}[ht]
    \centering
    \small
    \caption[Umfang und Typverteilung der englischen MultiNERD-Splits]{Umfang und Typverteilung der englischen MultiNERD-Splits.
    Tokens bezeichnen Wort-Tokens vor der Subwort-Tokenisierung;
    Entitäten wurden als Anzahl der \texttt{B}-Tags gezählt.}
    \label{tab:datensplits}
    \begin{tabular}{lrrr}
        \toprule
                                & \textbf{train} & \textbf{validation} & \textbf{test} \\
        \midrule
        Sätze                   & 131\,280    & 16\,410  & 16\,454 \\
        Tokens                  & 2\,842\,119 & 379\,011 & 347\,025 \\
        Entitäten               & 201\,769    & 26\,041  & 25\,645 \\
        \midrule
        PER (Person)            & 62\,987 & 7\,507 & 5\,265 \\
        ORG (Organisation)      & 27\,641 & 2\,737 & 3\,309 \\
        LOC (Ort)               & 58\,665 & 7\,850 & 12\,024 \\
        ANIM (Tier)             & 12\,736 & 1\,134 & 1\,604 \\
        BIO (biologische Entität) & 140   & 14     & 8 \\
        CEL (Himmelskörper)     & 2\,685  & 94     & 41 \\
        DIS (Krankheit)         & 8\,702  & 1\,695 & 759 \\
        EVE (Ereignis)          & 2\,525  & 299    & 352 \\
        FOOD (Lebensmittel)     & 8\,279  & 2\,120 & 566 \\
        INST (Instrument)       & 379     & 34     & 12 \\
        MEDIA (Medien, Werke)   & 6\,081  & 919    & 458 \\
        MYTH (mythologische Entität) & 569 & 56    & 32 \\
        PLANT (Pflanze)         & 7\,436  & 1\,180 & 894 \\
        TIME (Zeitangabe)       & 2\,540  & 322    & 289 \\
        VEHI (Fahrzeug)         & 404     & 80     & 32 \\
        \bottomrule
    \end{tabular}
\end{table}

Der \texttt{train}-Split dient vollständig der Optimierung der
Modellparameter bzw. Adapter, der \texttt{validation}-Split
ausschließlich der Checkpoint-Auswahl innerhalb eines Laufs
(Abschnitte~\ref{sec:encoder-pipeline} und~\ref{sec:llm-pipeline}); ein
systematischer Hyperparameter-Suchlauf findet nicht statt.
Der \texttt{test}-Split wird in der Hauptstudie je Lauf genau einmal
vollständig inferiert; ein separater Evaluationsschritt berechnet die
Metriken aus den gespeicherten Vorhersagen.

Einschränkend wurde der Testsplit bereits in den explorativen Vorläufen
ausgewertet; Tabelle~\ref{tab:datennutzung} ordnet jeder Entscheidung
Datengrundlage und Testbezug zu, einschließlich der nachträglichen
Sensitivitätsanalysen zum Ausgabelimit.
Für die trainierten Konfigurationen blieben die Hyperparameter zwischen
Vorlauf und Hauptstudie bis auf die Epochenzahl des 27B-Modells
unverändert. Die Läufe der Hauptstudie
wurden anschließend mit drei Seeds neu trainiert. Diese Wiederholungen
erfassen Trainingsvariabilität, stellen aber die Unabhängigkeit eines
bereits ausgewerteten Testsplits nicht wieder her.
Bei den Zero-Shot-Konfigurationen beruhen dagegen Decoding und
Prompt-Erweiterung auf Beobachtungen eines Vorlaufs auf dem Testsplit
(Abschnitt~\ref{sec:llm-pipeline}).
Die Zero-Shot-Ergebnisse werden deshalb als explorative Ergebnisse der
nach Testbeobachtungen gewählten Konfigurationen berichtet. Solche
Auswahlentscheidungen können zu einer optimistischen Verzerrung führen
\cite{cawley2010selectionbias}; deren Umfang lässt sich mit den vorhandenen Daten
nicht beziffern.

\begin{table}[ht]
  \centering
  \small
  \caption[Datengrundlage der Konfigurationsentscheidungen]{Datengrundlage
  der Konfigurationsentscheidungen und ihr Bezug zum MultiNERD-Testsplit.
  Direkt: Entscheidung anhand von Beobachtungen auf dem Testsplit;
  indirekt: im Vorlauf verwendet und nach dessen Testauswertung
  beibehalten; post hoc: nach Kenntnis der Testergebnisse ergänzt.}
  \label{tab:datennutzung}
  \begin{tabularx}{\linewidth}{@{}>{\raggedright\arraybackslash}p{3.0cm}%
      >{\raggedright\arraybackslash}X%
      >{\raggedright\arraybackslash}p{2.9cm}%
      >{\raggedright\arraybackslash}p{3.3cm}@{}}
    \toprule
    \textbf{Entscheidung} & \textbf{Datengrundlage} & \textbf{Testbezug} &
    \textbf{Folge} \\
    \midrule
    Thinking-Modus deaktiviert & generative Validierung eines
      abgebrochenen 4B-Vorlaufs (Parse-Fehlerraten 71 bis 83\,\%) &
      keiner & für alle Läufe übernommen \\
    \addlinespace
    Hyperparameter der trainierten Konfigurationen & im Vorlauf
      verwendet; nach dessen Testauswertung unverändert beibehalten &
      indirekt & keine Anpassung an Testwerte \\
    \addlinespace
    Epochenzahl 27B (2 auf 3) & Angleichung an die übrigen
      QLoRA-Konfigurationen; Testwert des Vorlaufs bekannt & indirekt &
      Hauptläufe mit drei Epochen \\
    \addlinespace
    Zero-Shot-Prompt (Glossar) und Sampling & Vorlauf von Qwen3.5-0.8B
      auf dem Testsplit (Greedy: fast nur leere Arrays) & direkt,
      leistungsbezogen & als explorative Ergebnisse berichtet;
      optimistische Verzerrung möglich \\
    \addlinespace
    Inferenzpräzision BF16 (Zero-Shot) & Durchsatzmessung auf einer
      Teilmenge des Testsplits (standardmäßig 96 Sätze) & direkt, nicht
      leistungsbezogen & BF16 statt NF4 des QLoRA-Regimes \\
    \addlinespace
    Ausgabelimit 256 Tokens & vor den Läufen ohne Längenanalyse
      festgelegt & keiner & Hauptauswertung bei 256 Tokens \\
    \addlinespace
    Analysen mit 1\,024 Tokens & Nachinferenz, Kontrollstichprobe und
      CrossNER-Neulauf nach Kenntnis
      der Testbefunde & post hoc & als Sensitivitätsanalysen berichtet \\
    \addlinespace
    Checkpoint-Auswahl & Validierungssplit (LLM: erste 200 Sätze) &
      keiner & -- \\
    \bottomrule
  \end{tabularx}
\end{table}

\section{Versuchsaufbau}
\label{sec:versuchsaufbau}

\subsection{Rechenumgebung und Reproduzierbarkeit}
\label{sec:rechenumgebung}

Alle Läufe wurden auf dem Slurm-Cluster des DFKI mit je einer GPU pro Job
ausgeführt.
Abhängig von Modellgröße und Verfügbarkeit der Partitionen kamen drei
GPU-Typen zum Einsatz (Tabelle~\ref{tab:hardware}); den angeforderten
Hauptspeicher nennt Tabelle~\ref{tab:app-einstellungen}.
Die MultiNERD-Läufe verwenden, soweit protokolliert, einheitliche
Softwareversionen (Anhang~\ref{app:durchfuehrung}).

\begin{table}[ht]
    \centering
    \small
    \caption[Eingesetzte GPU-Typen]{Eingesetzte GPU-Typen laut Laufmetadaten.
    Training und Inferenz einer MultiNERD-Konfiguration liefen jeweils
    auf demselben GPU-Typ; die Sensitivitätsanalysen mit 1\,024 Tokens
    (Abschnitte~\ref{sec:nachinferenz} und~\ref{sec:crossner}) liefen
    auf der H100, für das 4B- und das 27B-Modell also auf einem anderen
    GPU-Typ als ihre Hauptläufe.}
    \label{tab:hardware}
    \setlength{\tabcolsep}{4pt}
    \begin{tabular}{lll}
        \toprule
        \textbf{Modell} & \textbf{Regime} & \textbf{GPU} \\
        \midrule
        DeBERTa-v3-base, -large & Fine-Tuning & NVIDIA H100 (80\,GB HBM3) \\
    \addlinespace
        Qwen3.5-0.8B            & QLoRA       & NVIDIA H100 (80\,GB HBM3) \\
    \addlinespace
        Qwen3.5-4B              & QLoRA       & NVIDIA A100-SXM4 (80\,GB) \\
    \addlinespace
        Qwen3.5-27B             & QLoRA       & NVIDIA H100 NVL (94\,GB) \\
    \addlinespace
        Qwen3.5-0.8B, -4B       & Zero-Shot   & NVIDIA A100-SXM4 (80\,GB) \\
    \addlinespace
        Qwen3.5-27B             & Zero-Shot   & NVIDIA H100 NVL (94\,GB) \\
        \midrule
        DeBERTa-v3-large, Qwen3.5-27B & CrossNER-Inferenz & NVIDIA H100 (80\,GB HBM3) \\
    \addlinespace
        Qwen3.5-27B             & CrossNER (1\,024) & NVIDIA H100 (80\,GB HBM3) \\
    \addlinespace
        Qwen3.5, alle Größen    & Nachinferenz (1\,024) & NVIDIA H100 (80\,GB HBM3) \\
    \addlinespace
        Qwen3.5, alle Größen    & Kontrollstichprobe (1\,024) & NVIDIA H100 (80\,GB HBM3) \\
        \bottomrule
    \end{tabular}
\end{table}

Die Trainingsläufe der Hauptstudie laden Datensatz und Modelle über feste
Revisions-Hashes, setzen ihren Seed für Python, NumPy und PyTorch und
protokollieren je Lauf Git-Commit, Paketversionen, GPU- und Slurm-Angaben (Anhang~\ref{app:durchfuehrung}); die Zero-Shot-Läufe
luden Modelle und Datensatz dagegen ohne explizite Revisionsangabe und mit
reduzierten Metadaten.
Eine bitgenaue Reproduzierbarkeit wird für keinen der Läufe beansprucht;
die Zero-Shot-Ausgaben entstehen zudem durch Sampling, dessen
Generatorzustand bei einer Wiederaufnahme nicht wiederhergestellt wird.
Lauf, Commit und Ausgabeverzeichnis (Tabelle~\ref{tab:app-laeufe}),
vollständige Commits und Revisionen (Tabelle~\ref{tab:app-revisionen})
sowie nachrangige Einstellungen (Tabelle~\ref{tab:app-einstellungen})
stehen im Anhang.

\subsection{Encoder-basierte NER-Pipeline}
\label{sec:encoder-pipeline}

Die DeBERTa-v3-Checkpoints werden mit einem neu initialisierten
Klassifikationskopf aus Dropout (Rate 0,1) und linearer Schicht über 31
Klassen geladen und end-to-end auf dem vollständigen Trainingssplit
feinjustiert \cite{devlin2019bert,wolf2020transformers}.
Der Verlust ist die mittlere Token-Kreuzentropie über alle Positionen mit
Label (Abschnitt~\ref{sec:datenvorverarbeitung}), ohne Klassengewichte
oder Label Smoothing.
Optimiert wird mit AdamW (Standard des Trainers) und einem linearen
Lernratenplan mit Warm-up; Lernrate, Warm-up und Batchaufteilung je Modell
nennt Tabelle~\ref{tab:encoder-hyperparams}, die effektive Batchgröße
beträgt bei beiden Modellen 16.
Trainiert wird vollständig in FP32, da bei DeBERTa-v3 auf H100-GPUs mit
Mixed Precision numerische Instabilitäten (NaN-Werte) auftraten.

Nach jeder Epoche wird der vollständige Validierungssplit ausgewertet
(strikter, mikrogemittelter F$_1$-Wert;
Abschnitt~\ref{sec:performanzmetriken}) und ein Checkpoint gespeichert;
für die Testevaluation wird der Checkpoint mit dem höchsten
Validierungs-F$_1$ geladen.
Die konfigurierte Early-Stopping-Patience entspricht der maximalen
Epochenzahl, sodass das Kriterium nicht auslösen kann; alle sechs Läufe
durchliefen die volle Epochenzahl, die je Lauf gewählte Epoche nennt
Tabelle~\ref{tab:app-seeds} (Abschnitt~\ref{sec:erg-gesamt}).
In der Testinferenz wird jeder Satz einzeln verarbeitet und die Klasse
per argmax bestimmt.

\begin{table}[ht]
  \centering
  \small
  \caption[Hyperparameter der Encoder-Pipeline]{Hyperparameter der Encoder-Pipeline.
  Weitere Einstellungen (Adam-Parameter, Gradient Clipping,
  Checkpoint-Intervalle) enthält Tabelle~\ref{tab:app-einstellungen}.}
  \label{tab:encoder-hyperparams}
  \begin{tabular}{lll}
    \toprule
    \textbf{Hyperparameter} & \textbf{DeBERTa-v3-base} &
    \textbf{DeBERTa-v3-large} \\
    \midrule
    Lernrate & $2 \cdot 10^{-5}$ & $8 \cdot 10^{-6}$ \\
    Lernratenplan & linear, Warm-up 10\,\% & linear, Warm-up 6\,\% \\
    Mikro-Batch $\times$ Akkumulation & $16 \times 1$ & $8 \times 2$ \\
    Max.\ Epochen (Early-Stopping-Patience) & 5 (5) & 4 (4) \\
    Optimierer & \multicolumn{2}{l}{AdamW (fused), Weight Decay 0,01} \\
    Max.\ Sequenzlänge & \multicolumn{2}{l}{256 Subwort-Tokens} \\
    Numerische Präzision & \multicolumn{2}{l}{FP32} \\
    Seeds & \multicolumn{2}{l}{42, 123, 456} \\
    \bottomrule
  \end{tabular}
\end{table}

\subsection{LLM-basierte NER-Pipeline}
\label{sec:llm-pipeline}

\paragraph{Gemeinsame Einstellungen.}
Die Prompts werden mit dem Chat-Template des Tokenizers gerendert und um
den Beginn der Assistentenantwort ergänzt.
Der Thinking-Modus von Qwen3.5 ist in allen generativen Schritten
deaktiviert (Datengrundlage der Entscheidung in
Tabelle~\ref{tab:datennutzung} und Anhang~\ref{app:durchfuehrung}); das
Template setzt dazu einen leeren Denkblock vor die
Antwort, sodass das Modell unmittelbar mit der JSON-Ausgabe beginnt.
Die Antwort ist auf 256 neue Tokens begrenzt.
Dieser Wert wurde vor den Läufen ohne vorherige Längenanalyse
festgelegt; nachträglich zeigt sich, dass er die Referenzantworten des
Testsplits nahezu vollständig abdeckt: 55 der
16\,454 Referenzantworten überschreiten unter dem Qwen3.5-Tokenizer
256 Tokens und keine 1\,024.
Training und Validierung, an denen sich das Limit vorab hätte prüfen
lassen, zeigen dasselbe Bild: 1\,024 Tokens decken
alle Validierungsantworten und bis auf eine alle
Trainingsantworten ab (Längenverteilungen in
Anhang~\ref{app:durchfuehrung}).
Das Limit der Nachinferenz von 1\,024 Tokens
(Abschnitt~\ref{sec:nachinferenz}) ist damit ohne Rückgriff auf den
Testsplit begründbar.
Die Eingabe wird nicht gekürzt, und jeder Testsatz wird einzeln
verarbeitet.

\paragraph{Zero-Shot-Regime.}
Die Checkpoints werden unverändert in BF16 verwendet; alle drei Größen
erhalten denselben Prompt (Glossar-Profil,
Abschnitt~\ref{sec:prompt-design}) und dieselben Decoding-Parameter.
Abweichend vom QLoRA-Regime wird mit Sampling decodiert (Temperatur 1,0,
top-$k = 20$, top-$p = 1{,}0$, min-$p = 0$, Seed~42).
Die Einstellung übernimmt top-$k$ und min-$p$ der Qwen-Modellkarte
\cite{qwen2026card27b}; top-$p$ und Presence-Penalty weichen von deren
Empfehlungen für den Non-Thinking-Modus ab, die Temperatur entspricht
nur dem Profil für Reasoning-Aufgaben (Anhang~\ref{app:durchfuehrung}).
Diese Einstellungen wurden nach einem explorativen Vorlauf mit
Qwen3.5-0.8B festgelegt, in dem Greedy-Decoding für nahezu jeden Satz ein
leeres Array erzeugte (strikter F$_1$-Wert 0,5 Punkte; 16\,052 von 16\,454
Antworten waren leere Arrays).
Zugleich wurde der Prompt um ausgeschriebene Typbezeichnungen und eine
Definition benannter Entitäten erweitert, da die Abkürzungen des
Inventars ohne Training nicht selbsterklärend sind.
Decoding-Einstellung und Prompt-Änderung wurden damit gleichzeitig
eingeführt; ihre Einzelbeiträge sind nicht trennbar, und die Übertragung
der Sampling-Einstellung auf das 4B- und das 27B-Modell wurde nicht
eigens geprüft (Abschnitt~\ref{sec:limitationen}).
Für diese beiden Modelle erfolgte somit keine Konfigurationswahl anhand
eigener Testergebnisse.
Die Ergebnisse sind je eine Stichprobe aus der Ausgabeverteilung der
unangepassten Checkpoints unter genau diesen Prompt- und
Decoding-Bedingungen; ein isolierter Effekt der Anpassung ist wegen der
weiteren Änderungen gegenüber dem QLoRA-Regime nicht bestimmbar.

\paragraph{QLoRA-Regime.}
Die Anpassung erfolgt mit \texttt{peft} \cite{mangrulkar2022peft} und dem
\emph{SFTTrainer} aus \texttt{trl} \cite{vonwerra2026trldocs} als
Single-Turn-SFT: Jedes Trainingsbeispiel besteht aus Systemnachricht,
Nutzernachricht und Zielantwort im Standardprofil
(Abschnitt~\ref{sec:prompt-design}).
Multi-Turn-Ansätze wie ChatIE werden zugunsten eines einfachen, klar
evaluierbaren Setups nicht verwendet \cite{wei2023chatie}.
Entsprechend dem Standardverhalten von TRL~1.8.0 für Konversationsdaten
wird der Verlust über die gesamte Tokenfolge berechnet und nicht auf die
Assistentenantwort beschränkt (Option \texttt{assistant\_only\_loss},
Standardwert \texttt{False} \cite{vonwerra2026trldocs}); ein Teil des
Lernsignals entfällt damit auf die Reproduktion des konstanten Prompts.
Trainingssequenzen werden auf 1\,024 Tokens begrenzt und bei Überlänge
am Ende gekürzt, was auch die Zielantwort abschneiden kann; dies betrifft
48 der 131\,280 Trainings- und 11 der 16\,410 Validierungsbeispiele.
Sequence Packing wird nicht verwendet.

Zur Checkpoint-Auswahl erzeugt das Modell nach jeder Epoche per Greedy-Decoding
Vorhersagen für die ersten 200 Sätze des Validierungssplits, die wie in
der Testinferenz validiert und mit dem strikten F$_1$-Wert bewertet werden;
der zusätzlich auf dem vollständigen Validierungssplit protokollierte
Verlust dient nur der Überwachung.
Verwendet wird der Adapter der Epoche mit dem höchsten Wert, bei
Gleichstand der frühere; in acht der neun Läufe war dies die dritte, in
einem Lauf (4B, Seed~42) die zweite Epoche.
Die 200 Sätze weichen in den Typanteilen um höchstens 3,1 Prozentpunkte
vom Split ab, enthalten aber keine BIO- und INST-Entitäten
(Auszählung im Repository, Anhang~\ref{app:reproduzierbarkeit}).
Die nach F$_1$ gewählte Epoche stimmt nur in einem der neun LLM-Läufe
(4B, Seed~42) und in keinem der sechs Encoder-Läufe mit dem Minimum des
Validierungsverlusts überein.
Bei den LLMs liegt dieses Minimum in Epoche 1 oder 2.
Da 12 der 15 Läufe (8 von 9 LLM-, 4 von 6 Encoder-Läufen) die letzte
Epoche wählten, ist ein Untertraining nicht auszuschließen
(Abschnitt~\ref{sec:limitationen}).
Die Beschränkung auf 200 Sätze hält die generative Validierung je Epoche
im Bereich von Minuten (27B-Modell: rund 18 Minuten gegenüber rund
25 Stunden für den vollständigen Split), macht die Auswahl aber von einem
kleinen, nicht zufällig gezogenen Ausschnitt abhängig.
Ob eine größere oder zufällige Stichprobe dieselben Checkpoints gewählt
hätte, wurde nicht geprüft (Abschnitt~\ref{sec:limitationen}).
Für die Testinferenz wird der Adapter ohne Verschmelzung (Merge) auf das
4-Bit-quantisierte Basismodell aufgesetzt; die Ausgaben entstehen per
Greedy-Decoding.

\subsection{Prompt-Design und Ausgabeformat}
\label{sec:prompt-design}

Jeder Prompt besteht aus einer Systemnachricht mit Aufgabe, Typinventar
und Formatvorgaben sowie einer Nutzernachricht mit der nummerierten
Tokenliste im Format \texttt{Index: Token}.
Erwartet wird ein JSON-Array, dessen Objekte genau vier Felder
enthalten: \texttt{start\_token} (Index des ersten Tokens, inklusiv),
\texttt{end\_token} (Index nach dem letzten Token, exklusiv),
\texttt{text} (der abgedeckte Text) und \texttt{type} (eine der 15
Typabkürzungen); ohne Entitäten ist das leere Array zurückzugeben.
Eine Entität, die nur das Token mit Index~4 umfasst, wird also durch
\texttt{start\_token}~=~4 und \texttt{end\_token}~=~5 beschrieben.
Die Einordnung des Formats in die \emph{Indexing}-Familie nach
\textcite[Abschn.~3.2]{semin2026spanLabeling} gibt
Abschnitt~\ref{sec:k3-ausgabeformate}.

Es werden zwei Prompt-Profile verwendet, die jeweils wortgleich für alle
drei Modellgrößen gelten.
Das Standardprofil (\texttt{task\_v1}) nennt das Typinventar nur in Form
der Abkürzungen und wird für Training, generative Validierung und
Testinferenz im QLoRA-Regime verwendet (Beispielprompt und Zielantwort:
Listing~\ref{lst:prompt-task} in Anhang~\ref{app:durchfuehrung}).
Das Glossar-Profil (\texttt{multinerd\_\allowbreak glossary\_\allowbreak v2})
wird ausschließlich im Zero-Shot-Regime eingesetzt; es ergänzt die
Systemnachricht um eine Definition benannter Entitäten, die
ausgeschriebenen Typbezeichnungen und die Aufforderung, die gesamte
Tokenliste zu prüfen (Listing~\ref{lst:prompt-glossar} im Anhang).
Keines der Profile enthält Demonstrationsbeispiele.

Auf \emph{Constrained Decoding} wird bewusst verzichtet.
Verfahren der geführten Generierung begrenzen die Ausgabe auf die von
einer formalen Grammatik zugelassenen Formen
\cite{geng2023grammar}, können aber die Modellleistung beeinträchtigen
\cite{tam2024speakfreely}.
Freie Generierung mit nachgelagerter Validierung lässt dem Modell
dagegen das volle Ausdrucksspektrum, macht die Formtreue messbar und
entspricht Einsatzszenarien, in denen nur die Textausgabe einer API
zugänglich ist \cite{dagdelen2024structuredIE}.

\subsection{LoRA-Konfiguration und Hyperparameter}
\label{sec:lora-konfiguration}

LoRA wird gewählt, weil LLMs ohne aufgabenspezifisches Fine-Tuning bei
NER auf Tokenebene nur eingeschränkt konkurrenzfähig sind
\cite{lu2024llmstruggle}, Adapter den Speicher- und Rechenbedarf
gegenüber vollständigem Fine-Tuning erheblich senken
\cite{hu2021lora,han2024parameter,obeidat2025llm} und in den Vergleichen
von \textcite[Anh.~D, Tab.~11]{zhan2026generativeNER} zudem 0,14 bis 0,48
F$_1$-Punkte vor vollständigem Fine-Tuning liegen
(Abschnitte~\ref{sec:lora} und~\ref{sec:peft_ner_related}).
Der Rang beträgt $r = 16$ für die 0.8B- und 4B-Modelle und $r = 32$ für
das 27B-Modell, die Skalierung jeweils $\alpha = 2r$
(Tabelle~\ref{tab:lora-hyperparams}).
Bias-Terme werden nicht trainiert, Varianten wie rsLoRA oder DoRA nicht
eingesetzt.
Der höhere Rang des 27B-Modells folgt der Heuristik, größeren Modellen
mehr Adapterkapazität zuzuweisen; alternative Ränge wurden nicht
untersucht.

Die Adapter werden über die Namen von sieben Projektionstypen platziert:
\texttt{q\_proj}, \texttt{k\_proj}, \texttt{v\_proj} und
\texttt{o\_proj} in der Attention sowie \texttt{gate\_proj},
\texttt{up\_proj} und \texttt{down\_proj} im Feed-Forward-Netz.
Die Einbeziehung der Feed-Forward-Schichten folgt Befunden, nach denen
Adapter in unterschiedlichen Linearschichten komplementäre Kapazität
bereitstellen \cite{dettmers2023qlora,han2024parameter,wang2024parameter}.
Wegen der hybriden Architektur erfasst diese namensbasierte Auswahl
jedoch nicht alle Linearschichten: Die Attention-Projektionen existieren
nur in den Schichten mit voller Attention, während die anders benannten
Projektionen der Gated-DeltaNet-Schichten ohne Adapter bleiben.
Adaptiert werden somit die Feed-Forward-Netze aller Schichten und die
Attention jeder vierten Schicht (Modulzahlen in
Tabelle~\ref{tab:lora-hyperparams}).

Die Basismodelle werden mit \texttt{bitsandbytes}
\cite{dettmers2022llmint8} im 4-Bit-Format NormalFloat (NF4) geladen;
gerechnet wird in BF16, trainiert mit BF16-Mixed-Precision und Gradient
Checkpointing.
Nicht quantisierte Komponenten wie Embeddings und
Normalisierungsschichten bleiben in höherer Präzision.
Die gespeicherten Adapter der 0.8B- und 4B-Läufe liegen in BF16 vor, die
der 27B-Läufe nach der Wiederaufnahme aus Checkpoints
(Abschnitt~\ref{sec:effizienzmetriken}) in FP32
(Anhang~\ref{app:durchfuehrung}).
In zwei Punkten weicht die Umsetzung vom QLoRA-Standardrezept
\cite{dettmers2023qlora} ab: Double Quantization war vorgesehen, wurde
aber wegen eines fehlerhaft benannten Konfigurationsparameters nicht
aktiviert.
Nach \textcite[Abschn.~3]{dettmers2023qlora} spart sie im Mittel rund
0,37~Bit je Parameter ohne Leistungseinbuße, beim 27B-Modell also
höchstens rund 1,2~GiB des gemessenen Inferenzspeichers.
Zudem wird statt
eines Paged Optimizers der Standardoptimierer des Trainers verwendet
(AdamW, fusionierte PyTorch-Implementierung).
Die Konfiguration ist für alle drei Größen identisch, sodass sich
Größeneffekte nicht mit Unterschieden im Präzisionsregime vermischen.
Die Ergebnisse beziehen sich auf diese tatsächlich ausgeführte Variante.
Der Einfluss von Double Quantization, anderer Adapterplatzierung oder
eines auf die Antwort beschränkten Verlusts wurde nicht isoliert gemessen.

\begin{table}[ht]
  \centering
  \small
  \caption[LoRA-/QLoRA-Hyperparameter]{LoRA-/QLoRA-Hyperparameter.
  Der Anteil trainierbarer Parameter bezieht sich auf die Gesamtzahl aus
  Basismodell und Adaptern (0,759, 4,227 bzw. 27,055~Mrd.).
  Weitere Einstellungen enthält Tabelle~\ref{tab:app-einstellungen}.}
  \label{tab:lora-hyperparams}
  \begin{tabular}{lrrr}
    \toprule
    \textbf{Hyperparameter} & \textbf{0.8B} & \textbf{4B} & \textbf{27B} \\
    \midrule
    Rang $r$ / Skalierung $\alpha$ & 16 / 32 & 16 / 32 & 32 / 64 \\
    Adaptierte Module (FFN/Attention) & 72/24 & 96/32 & 192/64 \\
    Parameter Basismodell & 0,752\,B & 4,206\,B & 26,896\,B \\
    Trainierbare Parameter (Adapter) & 6,39\,M & 21,23\,M & 159,38\,M \\
    Anteil trainierbarer Parameter & 0,84\,\% & 0,50\,\% & 0,59\,\% \\
    Lernrate & $3 \cdot 10^{-4}$ & $2 \cdot 10^{-4}$ & $1 \cdot 10^{-4}$ \\
    Warm-up-Anteil (Kosinus-Schedule) & 5\,\% & 5\,\% & 3\,\% \\
    Mikro-Batch $\times$ Akkumulation & $8 \times 2$ & $4 \times 4$ & $1 \times 16$ \\
    \midrule
    \multicolumn{4}{@{}l}{\emph{Für alle Größen identisch:}} \\
    LoRA-Dropout & \multicolumn{3}{l}{0,05} \\
    Quantisierung / Rechenpräzision & \multicolumn{3}{l}{NF4 (ohne Double Quantization) / BF16} \\
    Optimierer & \multicolumn{3}{l}{AdamW (fused), Weight Decay 0,01} \\
    Effektive Batchgröße & \multicolumn{3}{l}{16} \\
    Epochen / Optimierungsschritte & \multicolumn{3}{l}{3 / 24\,615} \\
    Max.\ Sequenzlänge (Training) & \multicolumn{3}{l}{1\,024 Tokens} \\
    Seeds & \multicolumn{3}{l}{42, 123, 456} \\
    \bottomrule
  \end{tabular}
\end{table}

\subsection{Nachverarbeitung und Validierung der Ausgaben}
\label{sec:postprocessing}

Die generierten Antworten werden in beiden LLM-Regimen mit derselben
deterministischen, regelbasierten Kette in IOB2-Sequenzen überführt:

\begin{enumerate}
    \item \textbf{Bereinigung und Parsing:} Die Antwort wird von
    umgebendem Leerraum und Denkblöcken befreit und als JSON-Array
    geparst, bei Bedarf über zwei Rückfallstrategien
    (Anhang~\ref{app:durchfuehrung}).
    Liefert keine Strategie ein JSON-Array, gilt die Antwort als nicht
    parsebar.

    \item \textbf{Validierung:} Jedes Array-Element wird einzeln geprüft
    und verworfen, sobald eine der vier Regeln verletzt ist:
    \begin{itemize}
        \item \emph{Schema:} Das Element muss ein Objekt mit allen vier
        Feldern sein, wobei zusätzliche Felder zulässig sind und ignoriert
        werden; \texttt{text} und \texttt{type} müssen nicht-leere
        Zeichenketten sein.
        \item \emph{Indizes:} Start- und Endindex $s$ und $e$ müssen
        ganze Zahlen mit $0 \leq s < e \leq n$ sein ($n$: Tokenanzahl).
        \item \emph{Typ:} Der Typ muss nach Entfernen umgebenden
        Leerraums und ohne Beachtung der Groß- und Kleinschreibung im
        Inventar enthalten sein.
        \item \emph{Exakter Text:} \texttt{text} muss nach Entfernen
        umgebenden Leerraums exakt dem mit Leerzeichen verbundenen Text
        der indizierten Tokens entsprechen.
    \end{itemize}
    Gültige Typen werden auf die kanonische Schreibweise normalisiert;
    unbekannte Typen werden nicht auf Inventartypen abgebildet, sondern
    verworfen.

    \item \textbf{Konfliktauflösung:} Gültige Spans werden nach Start-,
    dann Endindex und schließlich Reihenfolge in der Antwort sortiert und
    nur übernommen, wenn sie keine bereits belegten Tokens überdecken.
    Damit setzt sich der früher beginnende, bei gleichem Beginn der
    kürzere Span durch; von exakten Duplikaten bleibt das erste erhalten.

    \item \textbf{IOB2-Konvertierung:} Ausgehend von einer reinen
    \texttt{O}-Sequenz erhält das erste Token jedes Spans das Label
    \texttt{B-}\emph{Typ}, jedes weitere \texttt{I-}\emph{Typ}.
\end{enumerate}

Nicht parsebare Antworten werden wie leere Vorhersagen behandelt und
ergeben eine reine \texttt{O}-Sequenz; bei parsebaren Antworten bleiben
gültige Elemente erhalten.
Verworfene Elemente zählen nicht als Vorhersage: Entsprach ein
verworfenes Element einer Referenzentität, bleibt diese unerkannt und
mindert den Recall; andernfalls entfällt lediglich eine falsch positive
Vorhersage.
Die Kette toleriert umgebenden Text, Markdown-Formatierung, Leerraum,
Typschreibweise und zusätzliche Felder, repariert aber keine Inhalte;
insbesondere wird syntaktisch fehlerhaftes JSON nicht instand gesetzt.

\section{Evaluationskonzept}
\label{sec:evaluationskonzept}

\subsection{Metriken der Erkennungsleistung}
\label{sec:performanzmetriken}

Zentrale Metriken sind strikte Precision, Recall und F$_1$ auf
Entitätsebene nach IOB2-Schema \cite{manning2008introduction,sang2003conll,ramshaw1995text}.
Eine Entität gilt nur als korrekt erkannt, wenn Spangrenzen und Typ exakt
mit der Referenz übereinstimmen, was dem \emph{strict}-Modus der
SemEval-2013-Auswertung entspricht \cite{segura2013semeval}.
Berechnet werden die Werte mit \texttt{seqeval}
\cite{nakayama2018seqeval} bei explizit gesetztem striktem Modus und
IOB2-Schema, sodass formal unzulässige Tag-Folgen keine Entität bilden
(Abschnitt~\ref{sec:evaluation}).
Primäre Kennzahl ist der mikrogemittelte F$_1$-Wert
\cite{manning2008introduction,keraghel2024recent}; fehlen Vorhersagen
oder Referenzentitäten, wird der betreffende Quotient als 0 definiert.
Beide Pipelines werden auf Wortebene gegen identische Referenzsequenzen
ausgewertet
(Abschnitt~\ref{sec:datenvorverarbeitung}).

Für trainierte Konfigurationen werden Mittelwert und
Stichprobenstandardabweichung über die drei Seeds berichtet, für die
Zero-Shot-Konfigurationen der jeweilige Einzellauf. Die
Standardabweichung beschreibt die beobachtete Streuung der Trainingsläufe
und dient nicht als Schwelle für die Bewertung von Unterschieden.
Die Unsicherheit durch die Auswahl der Testsätze wird mit einem
gepaarten Satz-Bootstrap beziffert, der dem Bootstrap nach
\textcite{efron1993bootstrap} folgt und wie bei
\textcite{koehn2004bootstrap} dieselben Satzstichproben für alle
verglichenen Systeme zieht.
Aus den 16\,454 Testsätzen werden 2\,000 Stichproben mit Zurücklegen
gezogen; für jede Stichprobe werden Mikro- und Makro-F$_1$ aller Läufe
sowie die Differenzen zwischen Konfigurationen neu berechnet.
Pro Stichprobe werden die Laufwerte über die drei festen Seeds
gemittelt; berichtet werden 95-\%-Perzentilintervalle (2,5- und
97,5-Perzentil) dieser Mittelwerte und ihrer Differenzen.
Bei Zero-Shot geht der jeweilige Einzellauf ein und wird gegen jeden der
drei Seeds der Vergleichskonfiguration gepaart.
Ein Differenzintervall, das null ausschließt, spricht unter diesen
Annahmen für eine gerichtete Differenz. Enthält es null, wird kein
eindeutiger Unterschied festgestellt; Gleichwertigkeit ist damit nicht
nachgewiesen \cite{lakens2017equivalence}. Typbezogene Vergleiche ohne eigene Intervalle werden
deskriptiv behandelt.

Das Bootstrap ist auf die vorhandenen trainierten Modelle und
gespeicherten Vorhersagen bedingt. Es erfasst weder die Unsicherheit
durch neue Trainingsläufe noch die Sampling-Streuung einer erneuten
Zero-Shot-Inferenz. Die Sätze werden als unabhängige Einheiten behandelt;
Abhängigkeiten innerhalb von Quelldokumenten werden nicht modelliert.
Als Hauptvergleiche gelten die aus FF1 und FF2 abgeleiteten
Gegenüberstellungen der Konfigurationen in
Tabelle~\ref{tab:experiment-matrix} bei 256 Ausgabetokens
(DeBERTa-v3-large gegen Qwen3.5-27B im QLoRA-Regime sowie die
Größenschritte innerhalb der Familien); die Vergleiche der nachträglichen Sensitivitätsanalysen
(1\,024 Tokens, CrossNER-Domänen) sind explorativ.
Die Intervalle sind nicht für Mehrfachvergleiche korrigiert. Sie beseitigen auch
keine mögliche Verzerrung durch die frühere Testnutzung.
Für die Fehleranalyse werden Precision, Recall, F$_1$ und Support zusätzlich
je Entitätstyp berechnet; sie sind für die Trainingsläufe der Hauptstudie
gespeichert und für alle Konfigurationen aus den Testvorhersagen
reproduzierbar.
Wegen der unausgewogenen Typverteilung (Abschnitt~\ref{sec:datenaufteilung},
Tabelle~\ref{tab:datensplits}) sind typbezogene Werte seltener Klassen mit
erheblicher Unsicherheit behaftet.

Ergänzend wird der makrogemittelte F$_1$-Wert als ungewichtetes Mittel der
typbezogenen F$_1$-Werte über alle 15 Typen berichtet
(Abschnitt~\ref{sec:evaluation}); er reagiert entsprechend stark auf die
seltenen Typen.
Für den Vergleich der Fehlerprofile wird jede Referenzentität genau einer
Klasse zugeordnet: korrekt erkannt (Grenzen und Typ stimmen überein),
\emph{Typfehler} (gleiche Grenzen, anderer Typ), \emph{Grenzfehler}
(überlappender Span, gleicher Typ), \emph{Grenz- und Typfehler}
(überlappender Span, anderer Typ) oder \emph{Auslassung} (keine
überlappende Vorhersage).
Vorhergesagte Spans ohne Überlappung mit einer Referenzentität zählen
gesondert als \emph{Zusätze}.
Die Zuordnung folgt der Logik der weniger strengen Evaluationsmodi von
\textcite{segura2013semeval}, dient aber nur der Beschreibung; bewertet
wird in allen Fällen strikt.

\subsection{Effizienzmetriken}
\label{sec:effizienzmetriken}

Für praxisnahe NER-Systeme ist neben der Erkennungsleistung die Effizienz
von Training und Inferenz zentral \cite{han2024parameter,obeidat2025llm}.

\paragraph{Trainingsdauer.}
Berichtet wird für alle Konfigurationen die Wall-Clock-Dauer der
Slurm-Trainingsjobs laut Abrechnung (\texttt{sacct}); sie umfasst Laden
von Modell und Daten, Training, Validierung nach jeder Epoche (bei den
LLMs inklusive generativer Validierung), Speichern von Checkpoints und
Export.
Die vom Trainer zusätzlich protokollierte Dauer ohne Laden und Export
liegt bei den ununterbrochenen Läufen darunter: bei den Encodern um 30
bis 49~s, bei den 0.8B- und 4B-Läufen um 32 bis 141~s (Maximum beim
4B-Modell mit Seed~42).
Die drei 27B-Läufe überschritten das Zeitlimit der Partition von drei
Tagen und wurden vom Scheduler dreimal neu eingereiht
(Anhang~\ref{app:durchfuehrung}); ihre Trainingsdauer ist die Summe der
vier Zuteilungen einschließlich des erneuten Ladens und mit den übrigen
Läufen nur näherungsweise vergleichbar.
Als Plausibilitätsprüfung dient die Hochrechnung der Schrittrate der
letzten Zuteilung auf alle Schritte.
Die Trainingsdauer hängt zudem von der Auslastung des Knotens ab: Die
drei 27B-Läufe liefen gleichzeitig auf demselben Knoten und benötigten
rund 38~s je Optimierungsschritt, der explorative Vorlauf mit zwei
Epochen lief allein auf dem Knoten und benötigte 15{,}5~s je Schritt.
Sie beschreibt daher den tatsächlich angefallenen Aufwand dieser Jobs;
ein von Auslastung und Neustarts unabhängiger Trainingsbedarf der
Modellkonfiguration wird daraus nicht abgeleitet.

\paragraph{Inferenzlatenz und Durchsatz.}
Gemessen wird je Testsatz bei Batchgröße~1 die Dauer des Modellaufrufs
(Vorwärtsdurchlauf bzw. Generierung), mit GPU-Synchronisation vor und
nach dem Aufruf und nach einem nicht gewerteten Aufwärmdurchlauf.
Nicht enthalten sind Laden, Tokenisierung, Datentransfer, Detokenisierung,
Validierung, Metrikberechnung und Speichern.
Berichtet werden Mittelwert, Median und 95.~Perzentil in Millisekunden
(für die Zero-Shot-Läufe wurde der Median nicht protokolliert),
für die LLMs zusätzlich der Durchsatz in Testsätzen je Sekunde reiner
Generierungszeit.
Die Werte beschreiben die Latenz einzelner Anfragen, nicht den mit
Batching erreichbaren Durchsatz; bei den LLMs hängen sie zudem von der
Länge der erzeugten Antwort ab.
Die Gesamtdauer der Testinferenz stammt aus der Phasenhistorie der
Laufverwaltung, bei den Zero-Shot-Läufen ohne Phasenhistorie aus der
mittleren Latenz hochgerechnet.

\paragraph{GPU-Speicher.}
Erfasst wird der maximal von PyTorch allokierte GPU-Speicher während der
Testinferenz in MiB; die Ergebnistabellen rechnen durch Division durch
1\,024 in GiB um.
Der Zähler wird nach dem Laden des Modells zurückgesetzt, sodass der Wert
die Modellgewichte und den zusätzlichen Inferenzbedarf umfasst;
reservierter Speicher und der über \texttt{nvidia-smi} sichtbare
Verbrauch werden nicht erfasst.
Der GPU-Speicherbedarf während des Trainings wurde nicht protokolliert.

\paragraph{Parameterzahlen.}
Berichtet werden Gesamtzahl und Zahl der trainierbaren Parameter
(Tabelle~\ref{tab:experiment-matrix}); ihr Verhältnis quantifiziert die Reduktion der aktualisierten Parameter durch
das Adaptertraining \cite{hu2021lora,dettmers2023qlora}.
Für die QLoRA-Modelle nennt Tabelle~\ref{tab:lora-hyperparams} Basismodell und Adapter getrennt,
da eine einfache Summierung der Tensorelemente die in 4~Bit gepackten
Gewichte nur zur Hälfte erfasst.

\paragraph{Vergleichbarkeit.}
Da die Konfigurationen auf unterschiedlichen GPU-Typen liefen
(Tabelle~\ref{tab:hardware}), enthalten Unterschiede in der Laufzeit
auch Hardwareeffekte (Beispiel in Anhang~\ref{app:durchfuehrung}).
Zero-Shot- und QLoRA-Inferenz unterscheiden sich zudem in Präzision und
Decoding.
Effizienzvergleiche werden daher vorrangig innerhalb gleicher Hardware-
und Präzisionsbedingungen interpretiert.

\paragraph{Ressourcenschwellen für FF3.}
Speicher- und Latenzbudgets wurden vor den Läufen nicht festgelegt.
FF3 wird deshalb über Schwellen beantwortet, die nachträglich aus den
gemessenen Werten abgeleitet sind: Mittlere Latenz und Inferenzspeicher
jeder Konfiguration geben an, ab welcher Grenze sie zulässig ist, und für
jede Kombination solcher Grenzen wird die zulässige Konfiguration mit dem
höchsten Mikro-F$_1$ bestimmt (Abschnitt~\ref{sec:disk-effizienz}).
Die Schwellen beschreiben die gemessenen Bedingungen (Batchgröße~1,
jeweiliger GPU-Typ), keine Anforderungen einer konkreten Anwendung.

\subsection{Formtreue der Ausgaben}
\label{sec:robustheit-ausgaben}

Während Encoder per Konstruktion genau ein Label pro Wort aus dem festen
Inventar liefern, können LLMs syntaktisch oder inhaltlich unzulässige
Ausgaben erzeugen.
Diese erfasst Tabelle~\ref{tab:formtreue} über Kennzahlen der
Validierungskette (Abschnitt~\ref{sec:postprocessing}) in vier
Fehlerklassen: \emph{JSON-Validität}
\cite{dagdelen2024structuredIE}, \emph{Schema-Treue},
\emph{Span-Alignment} \cite{semin2026spanLabeling} und
\emph{Konsistenz}.
Typen außerhalb des Inventars~\cite{tedeschi2022multinerd} gelten dabei
als halluzinierte Kategorien.
Für Encoder entfallen diese Kennzahlen; der Unterschied ist Teil des
Vergleichs, da LLMs Ausgabeformat und Typinventar erst über Prompt und
Fine-Tuning erlernen müssen \cite{dagdelen2024structuredIE,wei2023chatie}.

\begin{table}[ht]
  \centering
  \small
  \caption[Kennzahlen zur Formtreue der LLM-Ausgaben]{Kennzahlen zur Formtreue der LLM-Ausgaben in der Reihenfolge
  der Prüfung.
  Die Parse-Fehlerrate bezieht sich auf die Zahl der Testantworten, alle
  übrigen Werte sind absolute Häufigkeiten über den Testsplit.}
  \label{tab:formtreue}
  \begin{tabularx}{\linewidth}{@{}>{\raggedright\arraybackslash}p{3.3cm}%
      >{\raggedright\arraybackslash}X%
      >{\raggedright\arraybackslash}p{2.4cm}@{}}
    \toprule
    \textbf{Kennzahl} & \textbf{Definition} & \textbf{Fehlerklasse} \\
    \midrule
    Parse-Fehlerrate & Anteil der Antworten, aus denen keine Strategie ein
      JSON-Array gewinnt & JSON-Validität \\
    \addlinespace
    Parse-Strategien & Antworten, die direkt, nach Entfernen eines
      Markdown-Codeblocks bzw.\ per Extraktion aus umgebendem Text geparst
      wurden & JSON-Validität \\
    \addlinespace
    Falsche Struktur & gültiges JSON, das kein Array ist (je
      Parse-Versuch) & Schema-Treue \\
    \addlinespace
    Nicht-Objekte & Array-Elemente, die kein JSON-Objekt sind &
      Schema-Treue \\
    \addlinespace
    Fehlende Felder & fehlendes oder leeres Pflichtfeld & Schema-Treue \\
    \addlinespace
    Unzulässige Indizes & nicht ganzzahlig oder Verletzung von
      $0 \leq s < e \leq n$ & Span-Alignment \\
    \addlinespace
    Unbekannte Typen & Typ außerhalb des 15-Typen-Inventars &
      Schema-Treue \\
    \addlinespace
    Text-Abweichungen & \texttt{text} passt nicht zu den indizierten
      Tokens & Span-Alignment \\
    \addlinespace
    Überlappungen & wegen Überlappung mit einem übernommenen Span
      verworfen & Konsistenz \\
    \bottomrule
  \end{tabularx}
\end{table}

Bei der Interpretation sind drei Eigenschaften zu beachten.
Erstens sind die Fehlerzähler absolute Häufigkeiten; da alle
Konfigurationen denselben Testsplit verarbeiten, sind sie zwischen
Konfigurationen vergleichbar, erlauben aber keine elementbezogenen Raten,
weil die Gesamtzahl erzeugter Elemente nicht protokolliert wird.
Zweitens wird jedes verworfene Element nur der ersten verletzten Regel
zugeordnet, sodass die Prüfreihenfolge die Verteilung auf die Kategorien
bestimmt.
Drittens gilt eine Antwort bereits als geparst, wenn sie ein JSON-Array
enthält, auch wenn danach alle Elemente verworfen werden; die
Parse-Fehlerrate ist daher eine Untergrenze für den Anteil formal
fehlerhafter Antworten.
Ob eine Antwort am Ausgabelimit abbrach, wird nicht protokolliert
(Abschnitt~\ref{sec:nachinferenz}).

\subsection{Analyse des Einflusses der Parameteranzahl}
\label{sec:einfluss-parameteranzahl}

Der Einfluss der Parameteranzahl auf Erkennungsleistung, Effizienz und
Formtreue
wird getrennt nach Regime über die Größenklassen
(Abschnitt~\ref{sec:groessenklassen}) deskriptiv betrachtet.
Mit zwei Encoder- und drei LLM-Größen, bei denen Hyperparameter und
Hardware mitvariieren, lassen sich Größenunterschiede zwar gegen die
Variabilität zwischen den Seeds abgrenzen, Skalierungsgesetze im Sinne
angepasster Potenzgesetze \cite{kaplan2020scaling,hoffmann2022training} aber
nicht bestimmen.
Leitfrage ist das Verhältnis von Leistungsgewinn und Ressourcenbedarf
größerer LLMs gegenüber kompakten Encodern
\cite{keraghel2024recent,obeidat2025llm}.

\subsection{Nachinferenz am Ausgabelimit}
\label{sec:nachinferenz}

Das Limit von 256 Ausgabetokens ist eine Einstellung der Pipeline, die bei
langen Sätzen mit vielen Entitäten greift
(Abschnitt~\ref{sec:llm-pipeline}).
Um seinen Einfluss auf MultiNERD zu messen statt hochzurechnen, wurden für
alle neun QLoRA-Läufe genau die Testsätze, deren gespeicherte Antwort nicht
parsebar war, mit demselben Adapter, demselben Prompt und Greedy-Decoding,
aber einem Limit von 1\,024 Tokens erneut generiert; alle übrigen
Antworten wurden unverändert übernommen.
Die Auswahlregel \enquote{nicht parsebar} ist ein Stellvertreter für
\enquote{abgeschnitten}: Jede nicht parsebare Antwort der QLoRA-Läufe
endet ohne schließende Klammer, und eine vollständige, aber aus anderem
Grund nicht parsebare Antwort würde unter Greedy-Decoding identisch
reproduziert; abgeschnittene, aber zufällig parsebare Antworten lassen
sich in den gespeicherten Daten dagegen nicht erkennen.
Bei identischen Logits und Greedy-Decoding bliebe eine Antwort, die vor
dem Limit mit dem Ende-Token abschließt, unverändert; da die
ursprünglichen Läufe den Stoppgrund nicht protokollierten, ist diese
Bedingung für die übernommenen Antworten nicht vollständig prüfbar.
Zwei Kontrollen prüfen die Annahme.
Für jeden nachgenerierten Satz wird festgehalten, ob die ersten
256 erzeugten Tokens mit der gespeicherten Antwort übereinstimmen.
Zusätzlich wurden nachträglich je Lauf 1\,000 zufällig gezogene
Testsätze mit parsebarer gespeicherter Antwort mit demselben Adapter,
demselben Prompt und Greedy-Decoding bei einem Limit von 1\,024 Tokens
neu generiert und zeichenweise mit der gespeicherten Antwort verglichen
(\emph{Kontrollstichprobe}, Anhang~\ref{app:durchfuehrung}).
Nachinferenz und Kontrollstichprobe liefen für alle drei Größen auf der
H100, für das 4B- und das 27B-Modell also auf einem anderen GPU-Typ als
die ursprüngliche Testinferenz (Tabelle~\ref{tab:hardware}); da auch die
Encoder auf der H100 liefen, liefert die Kontrollstichprobe zugleich
Latenzen auf einheitlicher Hardware.
Die zusammengeführten Vorhersagen werden mit derselben strikten Metrik über
alle 16\,454 Sätze ausgewertet und je Seed den Werten bei 256 Tokens
gegenübergestellt.
Die Hauptauswertung beruht weiterhin auf dem Limit von 256 Tokens; die
zusammengeführte Nachinferenz wird als Sensitivitätsanalyse und nicht als
vollständiger Neulauf aller Testsätze berichtet. Beide Auswertungen gehen
in die Schlussfolgerungen ein
(Abschnitt~\ref{sec:erg-formtreue}).

\section{CrossNER-Ergänzungsversuch}
\label{sec:crossner}

Ergänzend wird untersucht, wie sich die auf MultiNERD angepassten Modelle
auf domänenspezifische Texte übertragen lassen.
Grundlage ist CrossNER mit den Domänen Künstliche Intelligenz, Literatur,
Musik, Politik und Naturwissenschaften und jeweils domänenspezifischen
Entitätstypen \cite{liu2021crossner}.
Die Daten stehen laut der Lizenzdatei des Repositorys
\texttt{zliucr/CrossNER} (Commit in Tabelle~\ref{tab:app-revisionen})
unter der MIT-Lizenz.
Der Versuch ist ein reiner Transfertest: Auf CrossNER findet weder
Training noch Checkpoint-Auswahl statt.

\paragraph{Daten und Typabbildung.}
Verwendet werden ausschließlich die Testsplits in einer festgelegten, über SHA-256-Prüfsummen verifizierten Revision: 431 (KI), 416
(Literatur), 465 (Musik), 651 (Politik) und 543 Sätze
(Naturwissenschaften), insgesamt 2\,506.
Ausgewertet werden nur die gemeinsamen Grobtypen PER, ORG und LOC, denen
die CrossNER-Typen gemäß Tabelle~\ref{tab:crossner-mapping} zugeordnet
werden; alle übrigen Referenztypen gelten als \texttt{O}.
Vorhersagen anderer MultiNERD-Typen werden ebenfalls auf \texttt{O}
abgebildet, Spangrenzen bleiben unverändert.
Vorhersagen der drei Grobtypen für Spans, die in CrossNER einem nicht
abgebildeten Typ wie \texttt{misc} angehören, zählen als falsch positiv.

\begin{table}[ht]
  \centering
  \small
  \caption[Abbildung der CrossNER-Typen auf die gemeinsamen Grobtypen]{Abbildung der CrossNER-Typen auf die gemeinsamen Grobtypen;
  alle übrigen CrossNER-Typen (z.\,B. \texttt{misc}, \texttt{event},
  \texttt{product}, \texttt{song}) werden als \texttt{O} behandelt.}
  \label{tab:crossner-mapping}
  \begin{tabularx}{\linewidth}{@{}lX@{}}
    \toprule
    \textbf{Grobtyp} & \textbf{CrossNER-Typen} \\
    \midrule
    PER & \texttt{person}, \texttt{politician}, \texttt{scientist},
      \texttt{writer}, \texttt{musicalartist}, \texttt{researcher} \\
    \addlinespace
    ORG & \texttt{organisation}, \texttt{politicalparty},
      \texttt{university}, \texttt{band} \\
    \addlinespace
    LOC & \texttt{location}, \texttt{country} \\
    \bottomrule
  \end{tabularx}
\end{table}

\paragraph{Modelle und Inferenz.}
Untersucht werden die jeweils größten Modelle beider Familien,
DeBERTa-v3-large und Qwen3.5-27B (QLoRA), mit den Checkpoints aller drei
Seeds; sechs eingefrorene Modelle auf fünf Domänen ergeben
30 Evaluationen.
Die Inferenz entspricht der MultiNERD-Testinferenz: Der Encoder
verarbeitet jeden Satz ungekürzt (Sätze über 256 Subwort-Tokens würden
zum Abbruch führen), das LLM erhält unverändert das Standardprofil mit
allen 15 Typen und erzeugt seine Ausgaben per Greedy-Decoding, ohne Thinking, mit
4-Bit-Basismodell und höchstens 256 Ausgabetokens.
Alle CrossNER-Evaluationen liefen auf NVIDIA H100 (80\,GB HBM3), für
Qwen3.5-27B also auf einem anderen GPU-Typ als die MultiNERD-Testinferenz
(Tabelle~\ref{tab:hardware}); Latenzwerte sind daher nur eingeschränkt
zwischen beiden Datensätzen vergleichbar.

\paragraph{Auswertung.}
Je Domäne werden strikte Precision, Recall und F$_1$ berechnet und je Seed
zu einem gepoolten Mikrowert über alle Domänen sowie zu einem gleich
gewichteten Domänenmittel aggregiert; berichtet werden Mittelwert und
Stichprobenstandardabweichung über die Seeds.
Für die Differenz zwischen Qwen3.5-27B und DeBERTa-v3-large wird zusätzlich
ein gepaartes Satz-Bootstrap mit 2\,000 Replikaten berechnet, das je
Domäne geschichtet so viele Sätze mit Zurücklegen zieht, wie der
jeweilige Testsplit enthält; dieselbe Ziehung gilt für beide Systeme,
alle Seeds und beide Ausgabelimits.
Die 95-\%-Perzentilintervalle unterliegen denselben Einschränkungen wie
das MultiNERD-Bootstrap (Abschnitt~\ref{sec:performanzmetriken}); die
domänenweisen Vergleiche sind explorativ und nicht für Mehrfachvergleiche
korrigiert.
Bezugsgröße ist der F$_1$-Wert derselben Modelle auf dem MultiNERD-Testsplit,
ebenfalls auf PER, ORG und LOC beschränkt.
Die Differenz $\Delta$F$_1$ beschreibt den beobachteten Unterschied zwischen
beiden Datensätzen; sie umfasst neben dem Domänenwechsel auch
Unterschiede in Annotationsrichtlinien, Typdefinitionen und
Textmerkmalen und ist daher nicht als isolierter Effekt des
Domänenwechsels zu interpretieren.
Zusätzlich werden für das LLM die Parse-Fehlerrate sowie für beide
Modelle Latenz und GPU-Spitzenspeicher erfasst.

\paragraph{Überschneidungen.}
21 CrossNER-Testsätze (KI: 2, Musik: 4, Politik: 3,
Naturwissenschaften: 12) stimmen in ihrer vollständigen Tokenfolge exakt
mit Trainingssätzen von MultiNERD überein; Überschneidungen mit
Validierungs- oder Testsplit bestehen nicht, auf Beinahe-Duplikate wurde
nicht geprüft.
Die Hauptauswertung umfasst alle 2\,506 Sätze; eine Sensitivitätsanalyse
wiederholt die Aggregation ohne die 21 Sätze.

\paragraph{Ausgabelimit.}
Die Wirkung des Limits von 256 Ausgabetokens (Abschnitt~\ref{sec:nachinferenz})
auf Recall und F$_1$ des LLM prüft Abschnitt~\ref{sec:erg-crossner}.
Eine Vorstudie auf den jeweils ersten 25 Sätzen der Domänen Musik und
Politik, einer nicht repräsentativen Stichprobe, gab den Anlass für
einen vollständigen Neulauf mit 1\,024 Ausgabetokens
(Anhang~\ref{app:durchfuehrung}).
Der Neulauf des 27B-Modells umfasst alle drei Seeds und alle 2\,506
Sätze der fünf Domänen; Prompt, Parser, Typabbildung und Greedy-Decoding
blieben unverändert, nur das Limit wurde angehoben.
Die Auswertung stellt beide Limits auf denselben Sätzen je Seed gepaart
gegenüber und wiederholt die Sensitivitätsanalyse ohne die
21 überschneidenden Sätze (Abschnitt~\ref{sec:erg-crossner},
Tabelle~\ref{tab:limit}).

\section{Validität und bekannte Abweichungen}
\label{sec:validitaet}

Tabelle~\ref{tab:abweichungen} fasst die beschriebenen Abweichungen vom
idealisierten Versuchsdesign mit Verweis auf die jeweiligen Abschnitte
zusammen; ihre Folgen für die Aussagekraft der Ergebnisse diskutiert
Abschnitt~\ref{sec:limitationen}.

\begin{table}[ht]
  \centering
  \small
  \caption[Bekannte Abweichungen vom idealisierten Versuchsdesign]{Bekannte Abweichungen vom idealisierten Versuchsdesign.}
  \label{tab:abweichungen}
  \begin{tabularx}{\linewidth}{@{}>{\raggedright\arraybackslash}p{3.3cm}%
      >{\raggedright\arraybackslash}X%
      >{\raggedright\arraybackslash}p{2cm}@{}}
    \toprule
    \textbf{Aspekt} & \textbf{Abweichung} & \textbf{Abschnitt} \\
    \midrule
    Referenz & Alle MultiNERD-Splits automatisch annotiert
      (Silberreferenz) & \ref{sec:datenauswahl} \\
    \addlinespace
    Hardware & Drei GPU-Typen; CrossNER-Inferenz von Qwen3.5-27B auf
      anderem GPU-Typ als MultiNERD & \ref{sec:rechenumgebung},
      \ref{sec:crossner} \\
    \addlinespace
    Zero-Shot-Setup & Glossar-Prompt, Sampling und BF16 statt
      Standardprofil, Greedy-Decoding und NF4; ein nicht
      deterministischer Lauf je Größe & \ref{sec:llm-pipeline} \\
    \addlinespace
    Testnutzung & Testsplit bereits in Vorläufen ausgewertet;
      Zero-Shot-Prompt und Decoding nach Testbeobachtungen gewählt;
      Präzision anhand von Durchsatzmessungen & \ref{sec:datenaufteilung},
      Tab.~\ref{tab:datennutzung} \\
    \addlinespace
    Herkunft & Zero-Shot-Läufe ohne explizite Revisionsangabe und mit reduzierten
      Metadaten; keine bitgenaue Reproduzierbarkeit &
      \ref{sec:rechenumgebung} \\
    \addlinespace
    LoRA-Platzierung & Keine Adapter in den Gated-DeltaNet-Schichten &
      \ref{sec:lora-konfiguration} \\
    \addlinespace
    QLoRA-Umsetzung & Ohne Double Quantization und Paged Optimizer &
      \ref{sec:lora-konfiguration} \\
    \addlinespace
    SFT-Verlust & Verlust über die gesamte Sequenz einschließlich Prompt
      & \ref{sec:llm-pipeline} \\
    \addlinespace
    Adapter-Präzision & Gespeicherte 27B-Adapter nach der Wiederaufnahme in
      FP32 statt BF16 & \ref{sec:lora-konfiguration} \\
    \addlinespace
    Hyperparameter & Größenabhängig gewählt, kein Suchlauf (Encoder und
      LLM) & \ref{sec:encoder-pipeline}, \ref{sec:lora-konfiguration} \\
    \addlinespace
    Checkpoint-Auswahl & LLM-Adapter nach den ersten 200
      Validierungssätzen, Encoder nach dem vollständigen Split &
      \ref{sec:llm-pipeline}, \ref{sec:encoder-pipeline} \\
    \addlinespace
    Längengrenzen & 256 Subwort-Tokens (Encoder), 1\,024 Tokens (SFT),
      256 Ausgabetokens (LLM-Inferenz, a priori ohne Längenanalyse
      gesetzt; Sensitivitätsanalysen mit 1\,024 Tokens) & \ref{sec:datenvorverarbeitung},
      \ref{sec:llm-pipeline}, \ref{sec:crossner} \\
    \addlinespace
    Effizienzmessung & 27B-Training in vier Slurm-Zuteilungen, Dauer aus
      der Slurm-Abrechnung; drei 27B-Seeds gleichzeitig auf einem Knoten;
      kein Trainingsspeicher protokolliert &
      \ref{sec:effizienzmetriken} \\
    \bottomrule
  \end{tabularx}
\end{table}
